\chapter{Conclusiones}
\label{cap:conclusiones}

Las conclusiones cierran la memoria y deben responder de forma directa al
problema planteado en la introducción. No conviene introducir aquí resultados
nuevos: lo adecuado es sintetizar lo demostrado en los capítulos anteriores.

\section{Conclusiones principales}
\label{sec:conclusiones-principales}

Resume las conclusiones más importantes del TFG en párrafos breves. Cada
conclusión debe estar apoyada por resultados, cálculos, pruebas o análisis ya
presentados en la memoria.

Un esquema posible es:

\begin{itemize}
    \item Se ha estudiado el contexto técnico y se han identificado los requisitos principales del problema.
    \item Se ha diseñado una propuesta coherente con las restricciones planteadas.
    \item Se han obtenido resultados que permiten valorar el comportamiento de la solución.
    \item Se han detectado limitaciones que deben tenerse en cuenta en futuras ampliaciones.
\end{itemize}

Sustituye esta lista por conclusiones específicas. Evita frases genéricas como
``se han cumplido todos los objetivos'' si no explicas cómo se ha comprobado.

\section{Cumplimiento de objetivos}
\label{sec:cumplimiento-objetivos}

Es recomendable relacionar cada objetivo específico de la
Sección~\ref{subsec:objetivos-especificos} con una evidencia concreta del
documento. La Tabla~\ref{tab:cumplimiento-objetivos} muestra un ejemplo.

\begin{table}[!ht]
    \footnotesize
    \centering
    \caption[Ejemplo de verificación del cumplimiento de objetivos]{Ejemplo de verificación del cumplimiento de objetivos. Fuente: elaboración propia}
    \label{tab:cumplimiento-objetivos}
    \begin{tabularx}{\textwidth}{p{4cm}X}
        \toprule
        \textbf{Objetivo} & \textbf{Evidencia en la memoria} \\
        \midrule
        Analizar el contexto & Capítulo~\ref{cap:marco-teorico} y revisión bibliográfica. \\
        Diseñar la propuesta & Capítulo~\ref{cap:metodologia}, requisitos y desarrollo. \\
        Evaluar resultados & Capítulo~\ref{cap:resultados}, pruebas y discusión. \\
        Identificar mejoras & Sección~\ref{sec:trabajo-futuro}. \\
        \bottomrule
    \end{tabularx}
\end{table}

\section{Limitaciones}
\label{sec:limitaciones-conclusiones}

Expón de forma honesta las limitaciones finales del trabajo. Algunas pueden
proceder de la metodología, otras de los recursos disponibles o del alcance
aprobado para el TFG.

Las limitaciones no deben redactarse como excusas, sino como condiciones que
ayudan a interpretar correctamente las conclusiones. Si una limitación afecta a
la validez de un resultado, indícalo expresamente.

\section{Trabajo futuro}
\label{sec:trabajo-futuro}

El trabajo futuro recoge mejoras razonables que podrían realizarse a partir del
TFG. No debe ser una lista de deseos inconexos, sino una continuación natural de
lo desarrollado.

Algunas líneas posibles son:

\begin{enumerate}
    \item Ampliar las pruebas con más datos, escenarios o condiciones de funcionamiento.
    \item Optimizar el diseño para reducir coste, tiempo de cálculo, consumo o mantenimiento.
    \item Integrar la propuesta en un entorno real o en una instalación piloto.
    \item Comparar la solución con nuevas alternativas técnicas.
    \item Desarrollar documentación, manuales o herramientas que faciliten su uso.
\end{enumerate}

Selecciona únicamente las líneas que tengan sentido para tu trabajo y explica
por qué serían valiosas.
